% ══════════════════════════════════════════════════════════════
%  SECTION 1: INTRODUCTION
% ══════════════════════════════════════════════════════════════

\section{Introduction}
\label{sec:intro}

Professional football matches nominally last 90~minutes, but the actual
duration has always been at the referee's discretion. Law~7 of the Laws
of the Game has required referees to add time lost to interruptions
since 1886, yet for decades the rule was systematically under-enforced:
added time averaged three to five minutes regardless of actual dead
time, which typically exceeded 28~minutes per match
\citep{forcher2024additional}. In November 2022, Pierluigi Collina,
chairman of FIFA's Referees Committee, instructed World Cup officials
to restore the actual time lost. The resulting matches---exceeding
100~minutes before a global audience of five billion---created a new
public benchmark for proper enforcement.

This paper studies what happened next, using data from 15~European
leagues and over 36,000~matches. We document three findings that
speak to broader questions in economics.

\textit{First}, the directive produced a dose-response gradient in
which enforcement intensity tracked institutional proximity to FIFA
with remarkable precision. The five major European leagues (the
``Big~Five'': England's Premier League, Spain's La~Liga, Germany's
Bundesliga, Italy's Serie~A, and France's Ligue~1) were formally
addressed, but the signal propagated far beyond them. The Premier
League shifted by 2.1~minutes (Cohen's $d = 1.13$); the Scottish
Premiership---nominally untreated---shifted by 1.5~minutes
($d = 0.83$); Serie~A, despite formal treatment, shifted by only
0.7~minutes ($d = 0.40$). This gradient reveals how a global
governing body's signal travels through the decentralised architecture
of European football, attenuating at each institutional boundary.

\textit{Second}, and most distinctively, the additional minutes
produced goal \textit{displacement} rather than goal creation. Total
goals per match were unchanged ($\hat{\beta} = -0.033$, $p = 0.39$),
a precisely estimated null. But goals migrated from the final minutes
of regulation into the expanded stoppage window: per-minute scoring
rates in the 76--84 minute window fell by 8.6\%, while the 90+ window
rose by 35.5\%. The share of matches in which the result changed
during stoppage time rose from 7.3\% to 10.3\%---roughly 55~additional
decisive moments per Big~Five season. We develop a simple model of
optimal attacking intensity under known remaining time to rationalise
this finding: when the trailing team expects more time, it substitutes
patience for desperation, spreading attacking effort across a longer
horizon without increasing aggregate intensity.

\textit{Third}, compliance was near-universal and immediate. Among
referees with sufficient pre- and post-directive observations,
95--100\% individually increased their stoppage time, with the
coefficient of variation across referees declining from 0.17 to 0.14.
The shift was permanent: three full seasons later, there has been no
reversion. The mechanism was focal-point coordination
\citep{schelling1960strategy}, not monitoring or sanctions---the World
Cup demonstration crystallised a new behavioural norm that became
self-reinforcing once a critical mass adopted it.

These findings contribute to several literatures. In sports economics,
the closest antecedents study referee bias under social pressure
\citep{garicano2005favoritism, dohmen2008professionals,
pricewolfers2010} and strategic manipulation of contests
\citep{dugganlevitt2002}. Our setting differs in that the rule was
unchanged---only its enforcement shifted---allowing us to study
compliance with an existing but under-enforced regulation, a problem
common across regulatory domains \citep{duflo2013truth,
johnson2020compliance}. The displacement result connects to the labour
economics literature on working hours and productivity
\citep{pencavel2015effects}: adding minutes to a football match is
analogous to extending a workday, and the finding of constant returns
to time has implications for how output responds to time extensions
when effort is endogenous. The compliance analysis speaks to the
literature on norm evolution and focal-point coordination
\citep{young1993evolution, rothockenfels2002}, showing how a single
high-visibility demonstration can resolve a coordination failure
among hundreds of independent professionals.

The paper proceeds as follows. Section~\ref{sec:first_stage} documents
the enforcement gradient across 15~leagues.
Section~\ref{sec:displacement} presents our central finding---goal
displacement---and develops a model of optimal attacking intensity that
rationalises it. Section~\ref{sec:data} describes the data.
Section~\ref{sec:strategy} presents the identification strategy,
including a dose-response framework that addresses the contamination
inherent in the binary difference-in-differences.
Section~\ref{sec:compliance} documents referee compliance and
convergence. Section~\ref{sec:conclusion} concludes.

